\documentclass[11pt]{article}
\usepackage[margin=1in]{geometry}
\usepackage{amsmath,amssymb,amsthm,mathtools,booktabs,array,longtable}
\usepackage[T1]{fontenc}
\usepackage{lmodern,microtype,xurl}
\usepackage[colorlinks=true,linkcolor=blue,citecolor=blue,urlcolor=blue]{hyperref}
\newtheorem{theorem}{Theorem}
\newtheorem{lemma}[theorem]{Lemma}
\newtheorem{proposition}[theorem]{Proposition}
\newtheorem{remark}[theorem]{Remark}
\newcommand{\D}{\mathcal D_5}
\newcommand{\W}{\mathcal W}
\newcommand{\Sym}{\mathfrak S}
\newcommand{\Q}{\mathbb Q}
\newcommand{\R}{\mathbb R}
\newcommand{\C}{\mathbb C}
\DeclareMathOperator{\diag}{diag}
\DeclareMathOperator{\sgn}{sgn}
\title{Sharp Five-Factor Recht--R\'e Lower Bounds\newline in Arbitrary Dimension}
\author{}
\date{Private AI-assisted draft, 19 October 2026\\Human review and priority assessment pending}
\begin{document}
\maketitle
\begin{abstract}
We determine a sharp lower Loewner bound for the average of ordered
products of five distinct positive semidefinite matrices, for every integer
$n\ge5$ and every finite matrix dimension, under $\sum_iA_i\preceq nI$.
Writing $r=n/(n-1)$ and
\[
 a=\frac{n\bigl(6n-9+\sqrt{20n^2-60n+49}\bigr)}{16(n-1)(n-2)},
\]
the optimal lower constant is
$\kappa_5(n)=a^2[r(3-r)a-r^2/2]$.
An exact uniform localizer certificate proves the universal inequality;
Zhong's known permutation-covariant real $n$-dimensional family attains it.
The proof reduces to four matrix polynomials of sizes $7,4,4,4$ and
104 positive semidefinite Bernstein coefficients over
$\Q(\sqrt5,\sqrt{249})$, together with two positive scalar blocks.
Separate exact implementations check the original free-word identity,
the complete all-$n$ representation map, and the positivity data.
This answers the five-factor specialization of Lai--Lim's question about
the correct lower bound. It does not settle arbitrary product lengths or
claim the first ordinary five-factor validity theorem.
The manuscript is private and unrefereed; independent AI-agent replay is
not human peer review or confirmed priority.
\end{abstract}

\section{Statement and relation to prior work}\label{sec:statement}
Let $(n)_5=n(n-1)(n-2)(n-3)(n-4)$ and define the ordered sum
\[
 \D(A)=\sum_{\substack{(i_1,\ldots,i_5)\in\{1,\ldots,n\}^5\\
                       i_1,\ldots,i_5\ \mathrm{distinct}}}
 A_{i_1}A_{i_2}A_{i_3}A_{i_4}A_{i_5}.
\]
Reversal of the five indices shows that this sum is Hermitian.
\begin{theorem}\label{thm:main}
For every integer $n\ge5$, let $A_1,\ldots,A_n$ be positive semidefinite
real symmetric or complex Hermitian matrices of any common finite dimension.
If $\sum_iA_i\preceq nI$, then
\begin{equation}\label{eq:main}
 \frac{\D(A)}{(n)_5}\succeq-\kappa_5(n)I,
 \qquad \kappa_5(n)=a^2\left[r(3-r)a-\frac{r^2}{2}\right],
\end{equation}
where
\begin{equation}\label{eq:parameters}
 r=\frac n{n-1},\qquad \Delta=\sqrt{20n^2-60n+49},\qquad
 a=\frac{n(6n-9+\Delta)}{16(n-1)(n-2)}.
\end{equation}
For every such $n$, a real $n$-dimensional tuple with $\sum_iA_i=nI$
attains equality on a nonzero vector. Thus $\kappa_5(n)$ is the exact
dimension-unrestricted lower constant over either field.
\end{theorem}
The universal part of the proof includes finite exact algebraic checks
specified in Sections~\ref{sec:localizer}--\ref{sec:verification}.
Section~\ref{sec:attainment} proves attainment directly. In raw-sum
normalization, set $C_5(n)=(n)_5\kappa_5(n)$; then
\[
 C_5(5)=\frac{3125(9033+581\sqrt{249})}{393216},\qquad
 \lim_{n\to\infty}\kappa_5(n)=\frac{11+5\sqrt5}{64}.
\]
The result is a lower endpoint, not a comparison with
$(n^{-1}\sum_iA_i)^5$, and not by itself an upper norm bound.
No minimum attaining dimension or equality classification is asserted.

Recht and R\'e proposed the noncommutative AM--GM problem~\cite{recht-re}.
Lai and Lim~\cite{lai-lim} disproved the original conjecture and developed
its noncommutative sum-of-squares formulation. Their Section~6 asks,
``What is the correct lower bound?'' Theorem~\ref{thm:main} answers that
question for product length five. Their numerical raw-sum value at
$(n,m)=(5,5)$ is prior evidence for this endpoint, not a new disproof here.
Zhong's Section~1.1 and Appendix~A~\cite{zhong} already supply the
permutation-covariant family and recurrence used below, building on
De Sa's examples~\cite{desa}. The new certificate matches the optimized
family value universally; the family and recurrence are not claimed as new.

The public candidate repository of \texttt{ipitchford}~\cite{pitchford}
already claims the ordinary five-factor upper bound for $n\ge5$ and lower
bound for $n\ge6$, with failure at $n=5$. Its word-orbit and symmetry-block
certificate architecture is reused and credited. Our target is the exact
lower constant rather than the unit bound. The inspected source is pinned
to commit \texttt{df3cb842014210eba99d1971be48aa079367feac}.
The literature audit is bounded: the full theorem of Liu's CISS 2025
paper~\cite{liu} was inaccessible through the checked lawful routes.
Its title and secondary abstract do not exclude overlap. Priority remains
unconfirmed, independently of the mathematical certificate checks.

\section{The known family attains the constant}\label{sec:attainment}
This section is a self-contained specialization of Zhong's family and
recurrence~\cite[Section~1.1 and Appendix~A]{zhong}.
In $\R^n$, put $e=n^{-1/2}(1,\ldots,1)^T$, $P=I-ee^T$, and
\[
 z_i=\sqrt r\left(e_i-\frac1n(1,\ldots,1)^T\right).
\]
Then $z_i\perp e$, $\|z_i\|=1$, $z_i^Tz_j=-1/(n-1)$ for $i\ne j$,
$\sum_i z_i=0$, and $\sum_i z_i z_i^T=rP$.
Let $d=n-1-(n-2)a$ and define
\begin{equation}\label{eq:family}
 A_i=a(P-z_i z_i^T)+(e+\sqrt d\,z_i)(e+\sqrt d\,z_i)^T.
\end{equation}
These are PSD once $a,d>0$. To check this for every $n\ge5$, write
$s=2-r\ge3/4$, $b=3-r$, and observe that $a$ is the larger root of
\begin{equation}\label{eq:stationarity}
 16s a^2-6rb a+r^2=0.
\end{equation}
For $1<r\le5/4$, this quadratic has two positive roots; its vertex
$3rb/(16s)$ is at most $5/8$. Its value at $1$ is
$7r^2-34r+32\ge7/16>0$. Therefore $0<a<1$, and
$d>n-1-(n-2)=1$. Summing~\eqref{eq:family} gives
\[
 \sum_i A_i=n ee^T+[na+r(d-a)]P=nI.
\]

For any ordered distinct-label word let
$t_k=e^TA_{i_k}\cdots A_{i_1}e$, with $t_0=t_1=1$.
Write the intermediate vector as $t_ke+\sum_{j=1}^k g_{j,k}z_{i_j}$
and $h_k=\sum_jg_{j,k}$. For a fresh label $i$,
\[
 A_i e=e+\sqrt d\,z_i,\qquad
 A_i z_j=az_j-\frac{\sqrt d}{n-1}e-\frac{d-a}{n-1}z_i\quad(j\ne i).
\]
Consequently
\[
 \binom{t_{k+1}}{h_{k+1}}=
 \begin{pmatrix}1&-\sqrt d/(n-1)\\
 \sqrt d&a-(d-a)/(n-1)\end{pmatrix}\binom{t_k}{h_k}.
\]
The transfer matrix has trace $2a$ and determinant $ar$. Thus
\begin{equation}\label{eq:recurrence}
 t_{k+2}=2at_{k+1}-ar t_k,
 \qquad
 t_5=8(2-r)a^4-4r(3-r)a^3+r^2a^2.
\end{equation}
Every distinct word has the same $e$-to-$e$ coefficient. Permutation
covariance makes $\D(A)e$ a multiple of $e$, so
$\D(A)e=(n)_5t_5e$. Using~\eqref{eq:stationarity},
\[
 t_5=a^2\left[\frac{r^2}{2}-r(3-r)a\right]=-\kappa_5(n).
\]
The stationary quadratic evaluated at $r/[2(3-r)]$ equals
$r^2[4(2-r)/(3-r)^2-2]<0$; hence the larger root exceeds that point
and $\kappa_5(n)>0$. This establishes the claimed attained value in
both the real and complex settings. No assumption that this family is
universally extremal is used: that conclusion follows from the certificate.

\section{A free-polynomial localizer certificate}\label{sec:localizer}
Use labels $0,\ldots,n-1$ in the certificate. Let $\W_n$ be the list
consisting of the empty word, the linear words, and all ordered quadratic
words, in that order and lexicographic order within each length.
Its real span has orthonormal word basis and dimension $1+n+n^2$.
For a word $w$, $w(A)^*$ reverses its letters; it does not commute them.

Let $G$ be a symmetric Gram matrix indexed by $\W_n$ and invariant
under permutations fixing label $0$. Define $G_i$ by relabelling $0$ as
$i$. Let $H$ be a symmetric Gram matrix invariant under all label
permutations. The certificate identity is
\begin{align}
 C_5(n)I+\D(A)
 &=\sum_i\sum_{w,v\in\W_n}(G_i)_{w,v}\,w(A)^*A_iv(A)\nonumber\\
 &\quad+\sum_{w,v\in\W_n}H_{w,v}\,
 w(A)^*\left(nI-\sum_iA_i\right)v(A).\label{eq:localizer}
\end{align}
If $G,H\succeq0$, factoring these scalar matrices expresses the right
side as sums of $p(A)^*A_ip(A)$ and
$q(A)^*(nI-\sum_iA_i)q(A)$. This proves positivity for the full domain
$\sum_iA_i\preceq nI$ in any finite dimension. The equality tuple is
used to choose certificate coordinates; the universal proof never replaces
the slack by zero.

\subsection{Canonical coordinates and the word identity}
For two words $w,v$, encode their concatenation with a separator between
them, canonically renaming new letters in order of first appearance.
For $G$, reserve $0$ for the marked label and name the others $1,2,\ldots$;
for $H$, start with $0$. Take the lexicographically smaller code obtained
from $(w,v)$ and $(v,w)$, with separator $-1$. Do not reverse either word
within this code. This describes exactly simultaneous label invariance
and Gram symmetry. There are $59$ marked and $22$ unmarked orbits.
Enumerating pairs in $\W_6\times\W_6$ and retaining each code on its
first occurrence fixes coordinates $z_0,\ldots,z_{80}$, with the marked
orbits first. At most four active labels occur in a pair, so this indexing
is valid for all $n\ge5$.

The exact rational functions $z_c(u)$ are the ordered list
\path{orbit_functions} in the frozen \path{CERTIFICATE.json};
Appendix~\ref{app:coordinates} specifies its 18 free coordinates and
the rational parameter $u$ is defined below. All $81$ functions are
machine-readable proof data, not numerical approximations.

To check~\eqref{eq:localizer}, take a concrete target word and sum over
its splits into a prefix, one inserted letter, and a suffix, with each
outer length at most two. The relevant $G_i$ entry has the reversed
prefix as its left word; subtract the corresponding $H$ entry. For
each split without an inserted letter, add $n$ times that $H$ entry.
The required coefficient is $C_5(n)$ for the empty word, $1$ for a
degree-five word with five distinct labels, and $0$ otherwise.
All words through degree five are covered by their $76$ equality
patterns. These coefficient equations have been checked as identities
in $\Q(u)$, not by evaluating finitely many $n$.
The target uses the raw ordered sum: there is no factor $5!$ or $(n)_5$
inside~\eqref{eq:localizer}. If $X_i=A_i/n$, the raw lower constant
becomes $C_5(n)/n^5$ on $\sum_iX_i\preceq I$.

\section{The complete invariant representation}\label{sec:representation}
We give the representation argument because checking a few contractions
of a Gram matrix is insufficient without completeness. Let $q$ be the
number of permuted letters: $q=n$ for $H$, and $q=n-1$ for $G$.
In the latter case $0$ is fixed and all sums in this section use the
other letters. In both cases $q\ge4$.

The invariant vectors for $H$, in order, are
\[
 1,\quad \sum_i x_i,\quad\sum_i x_i^2,\quad\sum_{i\ne j}x_ix_j.
\]
For $G$ they are
\[
 1,\ x_0,\ \sum_i x_i,\ x_0^2,\ \sum_i x_0x_i,\
 \sum_i x_ix_0,\ \sum_i x_i^2,\ \sum_{i\ne j}x_ix_j.
\]
These independent vectors have disjoint word supports and span the
entire fixed subspace. Contracting the Gram matrices against them defines
$H_0$ of size $4$ and $G_0$ of size $8$.

For $c\in\R^q$ with $\sum_i c_i=0$, define
\begin{align*}
 L(c)&=\sum_i c_ix_i,& Q(c)&=\sum_i c_ix_i^2,\\
 P(c)&=\sum_{i\ne j}(c_i+c_j)x_ix_j,&
 M(c)&=\sum_{i\ne j}(c_i-c_j)x_ix_j.
\end{align*}
These give four copies of the standard representation for $H$, in the
order $L,Q,P,M$. For $G$ use the six copies
\[
 L(c),\quad\sum_i c_ix_0x_i,\quad\sum_i c_ix_ix_0,\quad Q(c),P(c),M(c).
\]
They are orthogonal by support or by symmetry versus skew symmetry.
Their squared norm multipliers relative to $\|c\|^2$ are $1$ except
for $P,M$, where they are $2(q-2),2q$. Contract all copies against
the same $c=e_0-e_1$ for $H$, or $e_1-e_2$ for $G$, obtaining
$H_1$ of size $4$ and $G_1$ of size $6$.

The remaining pair spaces are the symmetric zero-diagonal arrays with
zero row sums, denoted $U_q$, and the skew arrays with zero row sums,
denoted $V_q$. Their dimensions are
\[
 \dim U_q=\frac{q(q-3)}2,\qquad
 \dim V_q=\frac{(q-1)(q-2)}2.
\]
Indeed, the symmetric row-sum map has rank $q$ for $q\ge3$ and the
skew row-sum map has rank $q-1$; their complementary subspaces are
exactly the pair copies already displayed.

For completeness, the permutation module on unordered pairs has a
three-dimensional commutant: an equivariant entry depends only on
intersection size $0,1,2$. It contains one trivial and one standard
summand, so its remaining nonzero summand $U_q$ has scalar commutant
and is absolutely irreducible. On skew pairs, the image of
$e_1\wedge e_2$ determines an equivariant map. Invariance under the
remaining permutations and antisymmetry under $(1\ 2)$ leave two
directions:
\[
 e_1\wedge e_2,\qquad
 \sum_{j>2}(e_1\wedge e_j-e_2\wedge e_j).
\]
Thus that commutant has dimension two. Removing its standard copy
leaves the absolutely irreducible $V_q$. The types are inequivalent
by dimension except at $q=4$, where the standard and $V_4$ both have
dimension three. Their transposition characters are $1$ and $-1$,
respectively, so they remain inequivalent. This covers the marked
$\Sym_4$ boundary at $n=5$.

Use the symmetric switch
\[
 x_ix_j+x_jx_i+x_kx_l+x_lx_k-x_ix_k-x_kx_i-x_jx_l-x_lx_j
\]
and the skew triangle
\[
 x_ix_j+x_jx_k+x_kx_i-x_jx_i-x_kx_j-x_ix_k
\]
as representatives of $U_q,V_q$. Their squared norms are $8,6$.
For $H$ take $(i,j,k,l)=(0,1,2,3)$; for $G$ take $(1,2,3,4)$.
The resulting scalar contractions are $H_2,H_3,G_2,G_3$.

\begin{lemma}\label{lem:blocks}
For every integer $n\ge5$, $G\succeq0$ if and only if
$G_0,G_1,G_2,G_3\succeq0$, and similarly for $H$.
\end{lemma}
\begin{proof}
The dimensions exhaust the word space:
\begin{align*}
 H:&\quad4+4(q-1)+q(q-3)/2+(q-1)(q-2)/2=1+n+n^2,\quad q=n,\\
 G:&\quad8+6(q-1)+q(q-3)/2+(q-1)(q-2)/2=1+n+n^2,\quad q=n-1.
\end{align*}
On each standard copy, equivariant bilinear forms are scalar multiples
of the standard inner product. Evaluation at the common $c$ multiplies
the multiplicity form by $\|c\|^2=2$, with invertible positive diagonal
changes for the embedding norms. On $U_q,V_q$, the contractions are
positive multiples $8,6$ of their scalar eigenvalues. Absolute
irreducibility and inequivalence leave no untested coupling.
The same real Gram matrices apply to complex Hermitian inputs.
\end{proof}
This also gives the symmetric invariant coordinate counts
$8\cdot9/2+6\cdot7/2+2=59$ and $4\cdot5/2+4\cdot5/2+2=22$.

\subsection{Exact construction of the block map}
For representatives $R_{B,i}$ just specified, each contracted entry is
\begin{equation}\label{eq:blockmap}
 B_{ij}(n)=\sum_{c=0}^{80}L_{B,ij,c}(n)z_c(u),\qquad
 L_{B,ij,c}(n)=\sum_{\substack{w,v\in\W_n\\\operatorname{orbit}(w,v)=c}}
 R_{B,i}(w)R_{B,j}(v),
\end{equation}
with the appropriate marked or unmarked coordinate offset. The left
word in this contraction is not reversed; reversal occurs only in
the free-polynomial coefficient extraction.
Fix the representative's anchor labels, including the mark for $G$.
There are respectively $1,3,5,4$ anchors for $G_0,G_1,G_2,G_3$, and
$0,2,4,3$ for $H_0,H_1,H_2,H_3$. A word pair contains at most four
fresh labels. A canonical fresh-label pattern with $k$ distinct labels
has $(n-\text{number of anchors})_k$ injective realizations.
This derives~\eqref{eq:blockmap} as a polynomial of degree at most four
in $n$. Impossible patterns automatically have falling factorial zero.

An independent implementation reconstructed the integer count maps at
$n=5,6,7,8,9$ and used exact degree-four interpolation. The preceding
degree bound makes this a proof of the count maps, not empirical
extrapolation of rational candidate entries. After substituting $n(u)$
and the $81$ functions $z_c$, all $81$ upper-triangular block entries
were checked as rational identities: $36+21+10+10+4=81$.
Thus the positivity data belong to the actual Gram matrices in
\eqref{eq:localizer}, rather than to an unrelated saved block array.

\section{The rational chart and kernel reductions}\label{sec:chart}
Put
\begin{equation}\label{eq:chart}
 E=16u^2-12u+1,\quad N=32u^2-18u+1,\quad
 n=\frac NE,\quad a=\frac N{16u-6},\quad
 u_0=\frac{3+\sqrt5}{8},\quad u_5=\frac{21+\sqrt{249}}{48}.
\end{equation}
On $u_0<u\le u_5$, $E>0$ and
\[
 n'(u)=-\frac{2(48u^2-16u+3)}{E^2}<0,
 \qquad n(u_5)=5,\qquad \lim_{u\downarrow u_0}n(u)=\infty.
\]
The numerator quadratic is positive by its negative discriminant.
The inverse branch is
\[
 u=\frac{6n-9+\Delta}{16(n-2)},\qquad a=\frac{nu}{n-1}.
\]
It satisfies~\eqref{eq:stationarity}, equivalently
$16(n-1)(n-2)a^2-6n(2n-3)a+n^2=0$.
Every finite integer $n\ge5$ therefore occurs exactly once.
Real-parameter polynomial positivity below is only a proof device;
the matrix-tuple theorem has integer $n$.

The denominator ledger is exhausted by the linear factors with roots
\[
 0,\ 1,\ 1/16,\ 1/2,\ 1/6,\ 1/8,\ 3/8
\]
and the three quadratics
\begin{equation}\label{eq:denoms}
 u^2-\frac{15u}{16}+\frac3{32},\quad
 u^2-\frac{17u}{16}+\frac9{64},\quad
 u^2-\frac{3u}{4}+\frac1{16}.
\end{equation}
The exact verifier checks divisibility into these factors with constant
remainder. They have no physical zero: $u_0>1/2$ and $u_5<37/48$;
the first two quadratic pairs are
$(15\pm\sqrt{129})/32$ and $(17\pm\sqrt{145})/32$, with small roots
below $1/2$ and large roots above $13/16$ and $29/32$. The final large
root is exactly the excluded $u_0$. These observations justify every
finite specialization, including $n=5$.

Let $\ell=n-(n-2)a$, $\beta=(n-2)a/(n-1)$, and
$\gamma=a[2(n-2)a-n]/(n-1)$. The exact blocks satisfy
\begin{align}
 G_0q_T&=0,& q_T&=(1,\ell,\beta,\ell^2,\beta\ell,\beta\ell,\beta\ell,\gamma)^T,
 \label{eq:kernelT}\\
 G_1q_S&=0,& q_S&=(1,a,a,\ell,a,0)^T,\nonumber\\
 G_1k&=0,& k&=(0,n-1,-(n-1),0,0,1)^T.\label{eq:kernelS}
\end{align}
The first two equations comprise the original fourteen kernel rows;
the last contributes six additional checked rows. Here
\[
 2(n-2)a-n=\frac{n(\Delta-2n-1)}{8(n-1)}>0,
 \quad \Delta^2-(2n+1)^2=16(n-1)(n-3)>0,
\]
so $\gamma>0$. Delete coordinate $7$ from $G_0$ (zero-based indexing).
For $G_1$, the rows $4,5$ of the two kernel columns $q_S,k$ form
$\diag(a,1)$; delete those coordinates.

\begin{lemma}\label{lem:delete}
If a Hermitian matrix $M$ annihilates the columns of $K$ and the rows
$K_J$ at a deleted coordinate set $J$ form an invertible square matrix,
then $M\succeq0$ if and only if the retained principal block is PSD.
\end{lemma}
\begin{proof}
For any vector $x$, let $y=x-KK_J^{-1}x_J$. It has zero coordinates
on $J$ and $x^*Mx=y^*My$. The converse is principal compression.
\end{proof}
Thus the remaining matrix sizes are $7,4,4,4$; no exact-rank assumption
at a special parameter is required. The scalar blocks, not imported
from a lower-degree certificate, are
\begin{align}
 G_2&=\frac{-1024u^2+1088u-144}{8u-3},&
 H_2&=\frac{-1024u^2+960u-96}{8u-3},\label{eq:scalars}\\
 G_3&=0,&H_3&=0.\nonumber
\end{align}
The first two are positive: their numerators are $-1024$ times the
second and first quadratics in~\eqref{eq:denoms}, respectively,
which are negative on the interval. Also $8u-3>0$ and $G_2-H_2=16$.

\section{The finite positivity certificate}\label{sec:verification}
Here $B$ denotes one of the retained blocks
$G_0[0{:}6],G_1[0{:}3],H_0,H_1$.
For the exponent vector $h_B$ in Table~\ref{tab:normalization}, set
\begin{equation}\label{eq:normalization}
 D_B(n)=\diag(n^{h_{B,i}-2}),\qquad
 P_B(u)=d_B(u)^2D_B(n(u))B(u)D_B(n(u)).
\end{equation}
The recorded general construction permits division by a common power
$E^{v_B}$ after squaring the denominator; the accepted data have
$v_B=0$ for all four blocks. Thus~\eqref{eq:normalization} is the exact
conversion used here. All entries of $D_B$ are nonzero and $d_B$ has
no physical zero, so it is an invertible congruence times a positive
scalar. It preserves PSD in both directions.

\begin{table}[ht]
\centering\small
\begin{tabular}{c c c c c}
\toprule
$B$ & size & $h_B$ & $\deg P_B$ & Bernstein degree $q_B$\\\midrule
$G_0$ &7&$(0,1,0,2,1,1,1)$&24&28\\
$G_1$ &4&$(1,2,2,2)$&12&16\\
$H_0$ &4&$(0,0,1,0)$&28&32\\
$H_1$ &4&$(1,2,1,1)$&20&24\\\bottomrule
\end{tabular}
\caption{Exact normalization and the four coefficient families.}\label{tab:normalization}
\end{table}
With $v=u-1/16$, $w=u-1/2$, $z=u-3/8$, the monic denominators are
\[
 d_{G_0}=uv^3w^3z^4,\quad d_{G_1}=vwz^3,\quad
 d_{H_0}=uv^4w^4z^4,\quad d_{H_1}=uv^3w^3z^2.
\]
The polynomial conversion is verified entry by entry as an exact
rational identity, including the indicated degrees.

Set $t=(u-u_0)/(u_5-u_0)$. Every $P_B$ has the exact expansion
\begin{equation}\label{eq:bernstein}
 P_B(u_0+(u_5-u_0)t)=
 \sum_{j=0}^{q_B}\binom{q_B}{j}t^j(1-t)^{q_B-j}C_{B,j},
 \qquad C_{B,j}\in\Q(\sqrt5,\sqrt{249})^{s_B\times s_B}.
\end{equation}
There are $29+17+33+25=104$ coefficient matrices. Each is symmetric
PSD and the four final coefficients $C_{B,q_B}$ are PD. The exact
matrices and their rational-field $LDL^T$ factors are stored in
\path{CERTIFICATE.json}. Its SHA-256 is
\begin{center}\footnotesize\ttfamily
 f4727ed67e787acb10208ca8fe7ea7a6c\\
 787cebfb117faf77df8647022843225.
\end{center}
This hash identifies the proof data; a hash alone is not a positivity
argument. The checks establishing their asserted properties are specified
next and have completed independent replays.

\subsection{Exact arithmetic and deterministic acceptance}
Represent each field element by four rational coordinates
\[
 A+B\sqrt5+C\sqrt{249}+D\sqrt{1245}
 =(A+B\sqrt5)+(C+D\sqrt5)\sqrt{249},
\]
with both radicals positive. Addition and multiplication follow their
quadratic relations. For $a+b\sqrt5$ with opposite nonzero signs,
the sign is that of $a$ if $a^2-5b^2>0$, and that of $b$ if it is
negative. Zero or matching signs are immediate. For
$x+y\sqrt{249}$ with $x,y\in\Q(\sqrt5)$ of opposite signs, compare
$x^2-249y^2$ in the smaller field in the same way. These norm tests
give exact real signs, with no tolerance or floating eigenvalue.
The rational basis is independent because $5$, $249$, and $1245$
are nonsquares, so coordinate equality is exact.

For each supplied factorization, the independent checker reconstructs
$C=LDL^T$, verifies symmetry and nonnegative diagonal entries of $D$,
and requires strictly positive diagonal entries for the final coefficient.
During factor construction, a zero pivot is allowed only with an entirely
zero remaining residual row. The replay can check the supplied factors
without repeating elimination. A factorization with a negative pivot
or failed identity is rejected.

To verify~\eqref{eq:bernstein}, first compose the rational polynomial
with $u_0+(u_5-u_0)t$ exactly. If its power coefficients are $p_i$, the
degree-$q$ Bernstein coefficients must satisfy
\[
 C_k=\sum_{i=0}^k\frac{\binom{k}{i}}{\binom{q}{i}}p_i,
 \qquad
 p_i=\binom qi\sum_{k=0}^i(-1)^{i-k}\binom ik C_k.
\]
Both are matrix identities. The independent replay used Horner
composition and forward differences, separate from the discovery
implementation, and checked every power coefficient including those
beyond the true degree. The verifier also reconstructs all $81$ orbit
functions, all $76$ word equations, all $20$ kernel rows, the four
scalar blocks, and all normalization and denominator identities.
The separate count-map replay described after~\eqref{eq:blockmap}
closes the remaining representation-map dependency.

\begin{proposition}\label{prop:certificate}
The frozen rational functions define PSD Gram matrices $G,H$ satisfying
\eqref{eq:localizer} for every integer $n\ge5$.
\end{proposition}
\begin{proof}
For $0<t\le1$, every scalar Bernstein weight in~\eqref{eq:bernstein}
is nonnegative, while the final weight is $t^{q_B}>0$.
Since $C_{B,q_B}\succ0$, all four $P_B(u)$ are PD on the physical
interval. Equation~\eqref{eq:normalization} gives PD of the retained
blocks. Lemma~\ref{lem:delete}, the verified kernels, and
\eqref{eq:scalars} give PSD of all eight full blocks. Their maps are
the contractions of the orbit Gram matrices, so
Lemma~\ref{lem:blocks} gives $G,H\succeq0$.
The exact $76$ coefficient identities give~\eqref{eq:localizer}.
No positivity assumption is needed at the excluded infinite-$n$
endpoint to infer positivity at finite $n$.
\end{proof}
Factoring the Gram matrices now proves the universal inequality in
Theorem~\ref{thm:main}; Section~\ref{sec:attainment} proves its sharpness.
Numerical optimization was used to discover a rational candidate, but
solver status, a floating margin, and rounding accuracy have no role in
this implication. Verification evaluates the fixed rational data only.

\section{Reproducibility, scope, and review status}\label{sec:record}
This \LaTeX{} source compiles without external inputs. The expanded
$81$ rational orbit functions, affine table, normalization polynomials,
and $104$ factorized coefficient matrices form a finite machine-readable
supplement. The companion private directory is
\path{overnight-ncagm-m5-all-n-private-package-20261019}; its
\path{README.md} specifies the data manifest. With CPython 3.12 and the
pinned SymPy 1.14.0 dependency already installed, its offline interface is
\begin{verbatim}
python3.12 -I verify.py
\end{verbatim}
No optimizer, network access, API credential, or original workspace is
required. The two stages use the byte-identical frozen certificate at
\path{stages/psd_words/INPUT-CERTIFICATE.json} and
\path{stages/block_map/CERTIFICATE.json}. The wrapper's only adaptation
of a mathematical worker removes three original absolute-environment
guard lines in its temporary copy; the exact diff is included.
A separate portable validation completed both stages with all exact checks
passing, no solver calls, and no mathematical core changes. Its receipt
is \path{portable.stdout.json} in the companion
\path{overnight-ncagm-m5-portable-validation-20261019} directory.
The source manifest beside this manuscript identifies that receipt and
the frozen research files replayed.

Two complementary independent exact implementations have completed:
the first checks the $81$ functions, $76$ original word patterns,
$20$ kernel rows, scalar conditions, polynomial transforms and all
$104$ coefficient factorizations; the second reconstructs the complete
all-$n$ block map and checks its $81$ upper-triangular identities.
They share exact rational-polynomial infrastructure but do not import
the discovery worker. The count-map replay uses exact integer counts
and the proved degree bound, not sampled positivity. Analytical reviews
separately establish representation completeness, physical-interval
signs, the operator bridge, and attainment. The acceptance record is
\path{ACCEPTANCE.json}, hash
\begin{center}\footnotesize\ttfamily
 c64dd723abacd28724f3a8472dee2df3f7\\
 caa725b91ceca04cf3d826d26084ff.
\end{center}

The proof and software were developed and audited by Codex AI agents
in a private campaign. ``Independent'' here means a separately
implemented check by another agent, not human refereeing or formal
verification in a proof assistant. No human author identity, credential,
endorsement, adoption, or employment consequence is asserted.
The manuscript has not been published or sent to the cited authors.
It settles only the stated five-factor lower-envelope problem; it does
not establish a universal extremizer for other product lengths, the
optimal general quadratic restoration threshold, or an algorithmic
convergence comparison. Broader priority checking and human mathematical
review remain open tasks.

\appendix
\section{Exact free-coordinate specification}\label{app:coordinates}
The $18$ free orbit columns $F_j$, their degrees $d_j$, and integers
$p_j,b_j$ in Table~\ref{tab:free} define
\begin{equation}\label{eq:free}
 z_{F_j}(n)=n^{4-d_j}\left[
 \frac{p_j}{2^{32}5^{4-d_j}}+
 \left(\frac1n-\frac15\right)\frac{b_j}{2^{32}}\right].
\end{equation}
Thus $z_{F_j}(5)=p_j/2^{32}$ is the fixed exact five-label seed.
These are the final accepted slopes, not the earlier rejected candidate.
The remaining coordinates are the exact rational strings in the
hash-bound \path{CERTIFICATE.json}; equivalently they are obtained
from \path{INPUT-AFFINE.json}'s $81\times19$
\path{solution_matrix_rows} applied to
$(z_{F_0},\ldots,z_{F_{17}},1)^T$.
Its hash is
\begin{center}\footnotesize\ttfamily
 7f8c881a0f58aa0d7e20f5e0b002b5c5\\
 d5b493299e8cb772d9eadd79f526a639.
\end{center}
The canonical indexing in Section~\ref{sec:localizer} fixes the meaning
of every column, so no unspecified ordering convention is required.
The validity check reconstructs the target identity from these functions;
it does not rely on the history of how the affine equations were solved.

\begin{table}[ht]
\centering\small
\begin{tabular}{r r r r r}
\toprule
$j$&$F_j$&$d_j$&$p_j=2^{32}z_{F_j}(5)$&$b_j$\\\midrule
0&17&3&-20817478263&88271623635\\
1&25&3&-18396979216&84921120947\\
2&61&2&11706295450&-9365392863\\
3&62&2&12981031495&-6837242709\\
4&64&2&75433334948&-81291345767\\
5&65&3&-24441626383&77610325069\\
6&66&3&-17077590016&71145486647\\
7&68&3&-12903045670&70652260805\\
8&69&3&-13157619467&57507367161\\
9&70&4&5969730599&-76348992586\\
10&71&4&2907939245&-55936656701\\
11&72&4&4181582164&-48912958583\\
12&73&4&2235060079&-67273566954\\
13&74&4&1929660160&-48211383911\\
14&76&4&3769223925&-104625015253\\
15&78&4&3242886733&-81769581127\\
16&79&4&5296781463&-98155140573\\
17&80&4&2541005637&-57879970852\\\bottomrule
\end{tabular}
\caption{Fixed dyadic seed and slope data in canonical orbit order.}\label{tab:free}
\end{table}

In JSON, a field coordinate is a four-element list of rational strings
in the order $(1,\sqrt5,\sqrt{249},\sqrt{1245})$.
The certificate's \path{coefficient_certificates} entries contain
\path{name}, \path{bernstein_degree}, and an ordered coefficient list.
Each list item has its index \path{k}, exact \path{matrix}, and
\path{LDL}, including the lower factor and diagonal pivots.
The \path{POLYNOMIAL-CHECKPOINT.json} entries specify $h_B,d_B,v_B$
and the polynomial affine rows before insertion of the fixed slopes.
All polynomial strings use the single variable \texttt{u} over $\Q$.
This specification, the listed finite data, and the exact arithmetic
rules in Section~\ref{sec:verification} determine a deterministic
certificate check without an optimizer or a choice of test points.

\begin{thebibliography}{9}
\bibitem{recht-re}
B. Recht and C. R\'e,
\emph{Toward a noncommutative arithmetic-geometric mean inequality:
conjectures, case-studies, and consequences}, COLT 2012, PMLR 23.
\href{https://proceedings.mlr.press/v23/recht12.html}{Primary proceedings page}.
\bibitem{lai-lim}
Z. Lai and L.-H. Lim,
\emph{Recht--R\'e Noncommutative Arithmetic--Geometric Mean Conjecture
is False}, ICML 2020, PMLR 119, 5608--5617.
\href{https://proceedings.mlr.press/v119/lai20a.html}{Primary proceedings page}.
\bibitem{desa}
C. De Sa, \emph{Random Reshuffling is Not Always Better}, NeurIPS 33 (2020).
\href{https://proceedings.neurips.cc/paper/2020/hash/42299f06ee419aa5d9d07798b56779e2-Abstract.html}{Primary proceedings page}.
\bibitem{zhong}
Y. Zhong, \emph{When the noncommutative AM-GM inequality holds},
arXiv:2610.04874v1, 4 October 2026.
\href{https://arxiv.org/abs/2610.04874v1}{Versioned primary record}.
\bibitem{pitchford}
\texttt{ipitchford}, \emph{Exact Low-Length Recht--R\'e Inequalities},
public candidate repository, \texttt{v1.1.1-candidate}.
Pinned commit \texttt{df3cb842014210eba99d1971be48aa079367feac}.
\href{https://github.com/ipitchford/exact-low-length-recht-re-inequalities/tree/df3cb842014210eba99d1971be48aa079367feac}{Versioned primary source}.
\bibitem{liu}
X. Liu, \emph{A Note on Arithmetic--Geometric Mean Inequality for
Well-Conditioned Matrices}, CISS 2025.
\href{https://doi.org/10.1109/CISS64860.2025.10944733}{DOI: 10.1109/CISS64860.2025.10944733}.
Full theorem not inspected; see the priority limitation in Section~\ref{sec:statement}.
\end{thebibliography}
\end{document}
